\chapter{The Theoretical Minimum}

We begin exactly where the lecture begins: by putting the spin system to one side for the moment and rebuilding the formalism already on the board. The mood is not that of a fresh textbook chapter starting from nowhere. It is a reset with intent. Susskind wants the earlier postulates back in view so that the next principle has somewhere natural to land. First come observables, eigenvalues, distinguishability, and probability. Then comes the new question, phrased just as one would phrase it classically: once the states are specified, how do they change with time? From there the lecture moves in a clean line toward unitary evolution, a Hermitian generator \(H\), the generalized time-dependent Schr\"odinger equation, and finally the evolution of averages, where the first strong analogy with classical mechanics appears.

\section{Recalling the Postulates of Measurement}

Susskind begins by saying that the postulates are not all independent, but that it is clearer to state them one at a time than to hunt for a minimal axiomatic set. He even remarks that observables are simply the things one measures, and that if the name had been up to him he would have called them \emph{measurables}. The standard language, however, is fixed, so we keep the usual term.

The lecture quickly rebuilds the measurement postulates in the following order:
\begin{enumerate}
\item Observables are represented by Hermitian operators.
\item The eigenvalues of such an operator are the possible outcomes of a measurement.
\item Physically distinguishable states are represented by orthogonal vectors.
\item Probabilities are determined by Born's rule.
\end{enumerate}

Along the way he inserts another statement which is not so much an additional postulate as a clarification of why Hermitian operators are the right objects: the result of a primitive experiment is a real number. If an apparatus seems to return a complex number, then in truth it is reporting two real quantities. Thus the eigenvalues that appear as measurement results must be real.

For the lecture, a generic observable is denoted by \(L\). Its eigenvectors and eigenvalues satisfy
\begin{equation}
L|\lambda\rangle=\lambda|\lambda\rangle.
\end{equation}
Here \(\lambda\) is a possible numerical outcome of measuring \(L\). Susskind is also willing to label the eigenvector by the same symbol \(\lambda\), and he openly admits that this is a notational shortcut. It is a useful one, and we will keep it.

There is one further point in the lecture that is worth preserving. Every observable is represented by a Hermitian operator, but the converse question, whether every Hermitian operator corresponds to a measurement one can actually perform in the laboratory, is a harder practical matter. The formal rule is liberal; the physics of building an apparatus is not.

The linear-algebra theorem that underlies the whole discussion is the familiar one:

\begin{theorem}
If \(L\) is Hermitian, then its eigenvalues are real. Eigenvectors with distinct eigenvalues are orthogonal, and within a degenerate eigenspace the eigenvectors may be chosen orthonormal. In particular, the eigenvectors of \(L\) can be chosen to form an orthonormal basis.
\end{theorem}

The lecture does not dwell on the proof here. It uses the theorem as the mathematical backbone for the measurement postulates.

\section{Distinguishability, Orthogonality, and Spin Examples}

The lecture next slows down and asks what \emph{physically distinguishable} really means. Rather than define it abstractly, Susskind gives it operational meaning.

Suppose someone hands us a spin and says: I prepared it either up or down. Can we tell which? Yes. Tilt the apparatus along the \(z\)-axis and measure. Up gives one answer, down gives the other, with no ambiguity. The same is true for left and right if we measure along the \(x\)-axis.

But now change the game. Suppose the preparation is either up or right. Then no single experiment, and no clever choice of orientation, can tell us unambiguously which preparation was used. If we measure along the \(z\)-axis, the up state gives \(+1\) with certainty, but the right state also has a nonzero probability of giving \(+1\). If we measure along the \(x\)-axis, the asymmetry reverses but the ambiguity does not disappear. This is the operational content of distinguishability.

\begin{definition}
Two states are \emph{physically distinguishable} if there exists an experiment, or a set of experiments, that can unambiguously determine which preparation was given.
\end{definition}

The lecture then promotes this operational notion into the vector-space postulate:

\begin{equation}
\text{physically distinguishable states} \quad \Longrightarrow \quad \text{orthogonal vectors.}
\end{equation}

This is the point at which the geometry of Hilbert space begins to carry physical meaning. Orthogonality is not just a formal relation; it is the mathematical translation of unambiguous distinguishability.

\subsection{Question \& Answer}

\paragraph{Question.} Can an experiment unambiguously distinguish a spin prepared up from a spin prepared right?

\paragraph{Answer.} No. Measuring along the \(z\)-axis will certainly identify the up state if the result is \(-1\), but the result \(+1\) is compatible with both preparations. Measuring along the \(x\)-axis produces the analogous problem in the opposite direction. The issue is not a poor choice of apparatus. The issue is that the two states are not completely distinguishable. In the Hilbert-space language, they are not orthogonal.

That is why the overlap between two states already contains physical information: it measures, in a precise sense, how far we are from complete distinguishability.

\section{Born's Rule and the Problem of Degeneracy}

Up to this point we know what values \emph{can} occur in a measurement. We do not yet know how likely they are. This is where, as Susskind says, the formalism acquires its real predictive bite.

Suppose the system is prepared in a normalized state \(|A\rangle\), and we measure the observable \(L\). If \(|\lambda\rangle\) is a normalized eigenvector of \(L\) with eigenvalue \(\lambda\), then the probability of obtaining that value is

\begin{equation}
P(\lambda)=|\langle \lambda|A\rangle|^2.
\end{equation}

The inner product \(\langle \lambda|A\rangle\) is the \emph{probability amplitude}. The lecture pauses over the obvious question: why the square? Because in a complex vector space the overlap itself need not be real, let alone positive. Its absolute square, on the other hand, is real and nonnegative, and so it can serve as a probability. For a normalized state the amplitudes also fit naturally with total probability:
\begin{equation}
\sum_i |\alpha_i|^2=1.
\end{equation}

If we expand the state in the eigenbasis of \(L\),
\begin{equation}
|A\rangle=\sum_i \alpha_i |\lambda_i\rangle,
\qquad
\alpha_i=\langle \lambda_i|A\rangle,
\end{equation}
then Born's rule becomes
\begin{equation}
P(\lambda_i)=|\alpha_i|^2.
\end{equation}

This is also the point where Susskind defends his notation. The label \(\lambda\) names an eigenvalue, but \(|\lambda\rangle\) is the eigenvector associated with it. A purist could insist on more indices, but the notation would become unreadable very quickly.

\subsection{Question \& Answer}

\paragraph{Question.} What if one eigenvalue corresponds to more than one eigenvector?

\paragraph{Answer.} Then the eigenvalue is degenerate, and the rule is still simple once we choose an orthonormal basis in that degenerate subspace: we square the overlaps for the orthogonal channels and add the resulting probabilities.

The lecture earns this answer rather than merely states it. If
\begin{align}
L|\lambda_1\rangle &= \lambda |\lambda_1\rangle,\\
L|\lambda_2\rangle &= \lambda |\lambda_2\rangle,
\end{align}
then any linear combination is again an eigenvector with the same eigenvalue:
\begin{equation}
L\bigl(\alpha|\lambda_1\rangle+\beta|\lambda_2\rangle\bigr)
=
\lambda\bigl(\alpha|\lambda_1\rangle+\beta|\lambda_2\rangle\bigr).
\end{equation}
So within a degenerate eigenspace we have freedom. In particular, we may choose an orthonormal basis there. Once that is done, the probability for the numerical outcome \(\lambda\) is the sum of the probabilities for the orthogonal channels:
\begin{equation}
P(\lambda)=\sum_a |\langle \lambda,a|A\rangle|^2.
\end{equation}
The extra label \(a\) is only a cleanup for the notes; the lecture itself keeps the discussion verbal. But the content is exactly his: first choose perpendicular vectors inside the degenerate subspace, then add the probabilities.

\section{Time Evolution, Reversibility, and the Unitary Operator}

Only after the kinematics of measurement are back in place does the lecture turn to a genuinely new question. In classical mechanics, once we know what the states are, the next question is how they move from one instant to the next. Susskind motivates the quantum version in explicitly classical terms.

For a classical coin, the law of motion could be trivial: put it down and nothing happens. Or it could be stroboscopic: at each tick it flips, heads to tails to heads to tails. What matters is that a law updates the state. Good classical laws are not only deterministic; they are reversible. Distinct states do not run into one another. If they did, we could not reconstruct the past from the present. Susskind calls this reversibility, information conservation, even the ``minus first law.''

Quantum mechanics must have its own version of that principle. The postulate is that the state at time \(t\) is obtained from the state at time \(0\) by the action of a single linear operator \(U(t)\), the same operator for every initial state:
\begin{align}
|\Psi(t)\rangle &= U(t)\,|\Psi(0)\rangle,\\
|\Phi(t)\rangle &= U(t)\,|\Phi(0)\rangle.
\end{align}
The lecture is careful on two points: \(U(t)\) is linear, and it does not depend on which state it acts on. It is the time-development operator for the system.

\begin{figure}[htbp]
\centering
\includegraphics[width=0.82\textwidth]{lecture_04_figure_02.png}
\caption{Time evolution under a common \(U(t)\) preserves orthogonality.}
\end{figure}

The blackboard argument is compact and worth preserving in clean form:
\begin{align}
|\Psi(t)\rangle &= U(t)\,|\Psi(0)\rangle, &
|\Phi(t)\rangle &= U(t)\,|\Phi(0)\rangle, \\
\langle \Psi(0)|\Phi(0)\rangle &= 0
&\Longrightarrow\quad
\langle \Psi(t)|\Phi(t)\rangle &= 0.
\end{align}
On the board Susskind is loose about capitals and lower-case letters; in the notes we standardize the second state to \(\Phi\).

This is the quantum form of information conservation: orthogonality breeds orthogonality. Physically distinguishable states evolve into physically distinguishable states.

The lecture then strengthens the statement. If orthogonality is preserved for a basis, then by expanding arbitrary states in that basis one shows that \emph{every} inner product is preserved. This is the ``nice little trick'' Susskind alludes to. Thus for arbitrary states,
\begin{equation}
\langle \Psi(t)|\Phi(t)\rangle
=
\langle \Psi(0)|\Phi(0)\rangle.
\end{equation}
The overlap remains constant. The degree of similarity between states does not change under the undisturbed time evolution.

\subsection{Question \& Answer}

\paragraph{Question.} Why does the bra version of the time-evolution equation involve \(U^\dagger(t)\)?

\paragraph{Answer.} Because a bra is the Hermitian conjugate of a ket. Starting from
\[
|\Psi(t)\rangle=U(t)|\Psi(0)\rangle,
\]
we take the Hermitian conjugate and obtain
\begin{equation}
\langle \Psi(t)|=\langle \Psi(0)|\,U^\dagger(t).
\end{equation}
The lecture later makes the same point for the Schr\"odinger equation itself: the derivative of a ket is a ket, the derivative of a bra is a bra, and passing from one to the other is Hermitian conjugation, not mere complex conjugation.

Now insert both the bra and ket evolutions into the overlap:
\begin{equation}
\langle \Psi(t)|\Phi(t)\rangle
=
\langle \Psi(0)|\,U^\dagger(t)U(t)\,|\Phi(0)\rangle.
\end{equation}
Since this must agree with \(\langle \Psi(0)|\Phi(0)\rangle\) for arbitrary initial states, the middle operator must act as the identity on every vector. Therefore
\begin{equation}
U^\dagger(t)U(t)=I.
\end{equation}
Such an operator is called \emph{unitary}. Equivalently,
\begin{equation}
U^\dagger(t)=U^{-1}(t).
\end{equation}

\begin{figure}[htbp]
\centering
\includegraphics[width=0.8\textwidth]{lecture_04_figure_03.png}
\caption{Unitarity and the infinitesimal evolution operator.}
\end{figure}

The board state that accompanies this step is slightly occluded, but its mathematical content is clear enough to reconstruct cautiously:
\begin{align}
U^\dagger(t)U(t)&=I,\\
U(\epsilon)&=I-i\epsilon H,\\
U^\dagger(\epsilon)&=I+i\epsilon H^\dagger.
\end{align}
The last line is a standard Hermitian-conjugate completion; the frame itself does not fully reveal the dagger on \(H^\dagger\).

\section{Infinitesimal Evolution, Hermitian \(H\), and Schr\"odinger's Equation}

The lecture now narrows the time step. Let \(\epsilon\) be a very small increment of time. Since \(U(0)=I\), the operator \(U(\epsilon)\) must lie very close to the identity. Susskind writes
\begin{equation}
U(\epsilon)=I-i\epsilon H
\end{equation}
to first order in \(\epsilon\). He pauses to note that the minus sign and the factor of \(i\) do not yet carry physical content. They are conveniences that can be absorbed into the definition of the operator \(H\). The real content is only that the deviation from the identity is of order \(\epsilon\).

Now impose unitarity on this infinitesimal form:
\begin{align}
I
&=
U^\dagger(\epsilon)U(\epsilon) \\
&=
\bigl(I+i\epsilon H^\dagger\bigr)\bigl(I-i\epsilon H\bigr) \\
&=
I+i\epsilon(H^\dagger-H)+O(\epsilon^2).
\end{align}
Ignoring terms of order \(\epsilon^2\), we find
\begin{equation}
H^\dagger=H.
\end{equation}
So the generator of infinitesimal time evolution must itself be Hermitian.

This is the moment at which the lecture christens \(H\). It is not \(H\) for \emph{Hermitian}; it is \(H\) for \emph{Hamiltonian}. Susskind is already willing to say that it is the energy, because it is Hermitian and because it plays in quantum mechanics the role the Hamiltonian plays in classical mechanics. But he is equally clear that the deeper identification will have to be checked against actual systems later in the course. At this stage we know the structural role of \(H\); we do not yet know which Hermitian operator nature chooses until we consult symmetry, prejudice, or experiment.

If the parameters of the problem do not depend explicitly on time, then \(H\) does not depend explicitly on time either. This is the time-translation-invariant case. If the external conditions are varied in time, \(H\) can depend on time, just as in classical mechanics. The lecture already warns us that in that case simple energy conservation is lost.

With the infinitesimal form of \(U\) in hand, the state update becomes
\begin{equation}
|\Psi(\epsilon)\rangle=(I-iH\epsilon)\,|\Psi(0)\rangle.
\end{equation}

\begin{figure}[htbp]
\centering
\includegraphics[width=0.82\textwidth]{lecture_04_figure_04.png}
\caption{From infinitesimal evolution to a difference quotient.}
\end{figure}

The corresponding blackboard algebra is
\begin{align}
|\Psi(\epsilon)\rangle &= (1-iH\epsilon)\,|\Psi(0)\rangle,\\
\frac{1}{\epsilon}\Bigl(|\Psi(\epsilon)\rangle-|\Psi(0)\rangle\Bigr)
&=
-iH|\Psi(0)\rangle.
\end{align}
Only the beginning of the derivative line is visible in this frame, so we do not claim more from the image than it actually shows.

Now comes the key conceptual step. The initial time need not literally be \(0\); it can be any time, so long as \(\epsilon\) is measured from it. The left-hand side is therefore the usual difference quotient for a vector-valued function. Susskind explicitly remarks that there is no mystery here: if one can subtract vectors, then one can divide the difference by a number, and in the limit one gets a derivative. Thus
\begin{equation}
\frac{d}{dt}|\Psi(t)\rangle=-iH|\Psi(t)\rangle.
\end{equation}

\begin{figure}[htbp]
\centering
\includegraphics[width=0.82\textwidth]{lecture_04_figure_05.png}
\caption{The boxed time-dependent Schr\"odinger equation on the board.}
\end{figure}

The boxed board result is the one the lecture is heading toward, and it is worth echoing typographically:
\begin{equation}
\fbox{$\displaystyle \frac{d}{dt}|\Psi(t)\rangle=-iH|\Psi(t)\rangle$}
\end{equation}
This is the generalized time-dependent Schr\"odinger equation in units where \(\hbar=1\).

Later in the lecture Susskind restores dimensions by recalling that \(H\) has units of energy while the time derivative has units of inverse time. To repair the dimensions one inserts Planck's constant:
\begin{equation}
i\hbar\,\frac{d}{dt}|\Psi(t)\rangle=H|\Psi(t)\rangle.
\end{equation}

\subsection{Question \& Answer}

\paragraph{Question.} How can this reversible, unitary evolution be correct if a measurement can disturb a state and send it into something apparently different?

\paragraph{Answer.} Because the equation we have derived governs an \emph{isolated} system. If a measuring apparatus interacts strongly enough with the system to matter, then the apparatus must be included as part of a larger quantum system. One cannot disturb the system and then ignore the degrees of freedom that did the disturbing. The lecture therefore postpones the real discussion of measurement until entanglement has been introduced.

\section{Expectation Values and the Evolution of Averages}

To explain what the new equation has to do with classical mechanics, the lecture changes focus. It no longer asks directly about the state vector; it asks about averages of observable quantities.

Susskind immediately complains about the standard name. The ``expectation value'' is really an average value, and averages can be misleading if one forgets what they are. A fair coin assigned the values \(+1\) and \(-1\) has average \(0\), even though \(0\) is never observed. A probability distribution can even have zero probability in a neighborhood of the average. So the phrase ``expected value'' is, if anything, even worse. Still, the standard terminology is entrenched, and we keep it.

If an observable can take values \(\lambda_i\) with probabilities \(P(\lambda_i)\), the average is defined by
\begin{equation}
\bar{\lambda}=\sum_i \lambda_i\,P(\lambda_i).
\end{equation}

Now let \(L\) be a Hermitian operator with orthonormal eigenbasis \(\{|\lambda_i\rangle\}\), and expand the state as
\begin{equation}
|A\rangle=\sum_i \alpha_i |\lambda_i\rangle.
\end{equation}
By Born's rule,
\begin{equation}
P(\lambda_i)=|\alpha_i|^2=\alpha_i^*\alpha_i.
\end{equation}
Therefore the average value of the measured outcome is
\begin{equation}
\bar{\lambda}=\sum_i \lambda_i |\alpha_i|^2.
\end{equation}

The lecture then proves that this is the same as sandwiching the operator \(L\) between the bra and ket of the state.

\begin{proof}
The bra corresponding to \(|A\rangle\) is
\begin{equation}
\langle A|=\sum_j \alpha_j^* \langle \lambda_j|.
\end{equation}
Hence
\begin{align}
\langle A|L|A\rangle
&=
\sum_{i,j}\alpha_j^*\alpha_i\langle \lambda_j|L|\lambda_i\rangle \\
&=
\sum_{i,j}\alpha_j^*\alpha_i\,\lambda_i\,\langle \lambda_j|\lambda_i\rangle \\
&=
\sum_i \lambda_i |\alpha_i|^2,
\end{align}
because the orthonormality relation \(\langle \lambda_j|\lambda_i\rangle=\delta_{ji}\) collapses the double sum to a single sum.
\end{proof}

Thus
\begin{equation}
\langle L\rangle=\langle A|L|A\rangle.
\end{equation}
This is where the bracket notation earns its keep: the average is a bra-ket sandwich.

The lecture is happy to use bracket notation and bar notation interchangeably. If we want to emphasize time dependence, we may write
\begin{equation}
\bar L(t)=\langle \psi(t)|L|\psi(t)\rangle.
\end{equation}
The time dependence is in the state vectors. The operator \(L\) itself is fixed.

Now differentiate:
\begin{equation}
\frac{d}{dt}\langle L\rangle
=
\left(\frac{d}{dt}\langle \psi|\right)L|\psi\rangle
+
\langle \psi|L\left(\frac{d}{dt}|\psi\rangle\right).
\end{equation}
In units where \(\hbar=1\), the ket equation is
\begin{equation}
\frac{d}{dt}|\psi\rangle=-iH|\psi\rangle,
\end{equation}
and the bra equation, obtained by Hermitian conjugation, is
\begin{equation}
\frac{d}{dt}\langle \psi|=i\langle \psi|H.
\end{equation}
Substituting both into the derivative gives
\begin{align}
\frac{d}{dt}\langle L\rangle
&=
i\langle \psi|HL|\psi\rangle
-
i\langle \psi|LH|\psi\rangle \\
&=
i\langle \psi|(HL-LH)|\psi\rangle.
\end{align}
This naturally introduces the commutator,
\begin{equation}
[H,L]=HL-LH,
\end{equation}
and so
\begin{equation}
\frac{d}{dt}\langle L\rangle=i\langle [H,L]\rangle.
\end{equation}
Restoring \(\hbar\) gives the dimensionally correct form
\begin{equation}
\frac{d}{dt}\langle L\rangle=\frac{i}{\hbar}\langle [H,L]\rangle.
\end{equation}
Because commutators are antisymmetric, this may also be written as
\begin{equation}
\frac{d}{dt}\langle L\rangle=-\frac{i}{\hbar}\langle [L,H]\rangle.
\end{equation}
This second form is the one that will line up most directly with the classical Poisson-bracket formula.

\section{Classical Correspondence, Poisson Brackets, and Energy Conservation}

Before finishing, the lecture deliberately glances back at classical mechanics. The goal is not to found quantum mechanics on classical intuition, but to see how classical-looking equations emerge from the quantum formalism.

In classical Hamiltonian mechanics, the time evolution of a quantity \(L\) is written
\begin{equation}
\dot L=\{L,H\}_{\mathrm{PB}},
\end{equation}
where \(\{\cdot,\cdot\}_{\mathrm{PB}}\) denotes the Poisson bracket. Compare this with
\begin{equation}
\frac{d}{dt}\langle L\rangle=-\frac{i}{\hbar}\langle [L,H]\rangle.
\end{equation}
The similarity is too striking for the lecture to ignore.

The parallel deepens because Poisson brackets and commutators share the same formal properties. Both are antisymmetric:
\begin{align}
\{L,H\}_{\mathrm{PB}}&=-\{H,L\}_{\mathrm{PB}},\\
[L,H]&=-[H,L].
\end{align}
Both also satisfy the same product rule pattern:
\begin{align}
\{LM,H\}_{\mathrm{PB}}&=\{L,H\}_{\mathrm{PB}}\,M+L\,\{M,H\}_{\mathrm{PB}},\\
[LM,H]&=L[M,H]+[L,H]M.
\end{align}
The lecture explicitly points out the second identity as a one-line exercise in adding and subtracting \(LHM\).

At this stage the connection is still formal, and Susskind is careful about the philosophical order. Dirac recognized that commutators and Poisson brackets ought to be related, and historically that observation was enormously fruitful. But logically the lecture insists on the opposite direction: quantum mechanics comes first, classical mechanics later, as an approximation. The notes should keep that tone. What we have here is not a derivation from classical mechanics, but a clue about how the classical shadow must look.

With the commutator convention \([A,B]=AB-BA\), the correspondence suggested by the lecture may be written in either of the equivalent forms
\begin{equation}
\{A,B\}_{\mathrm{PB}}
\;\longleftrightarrow\;
-\frac{i}{\hbar}[A,B],
\qquad\text{or}\qquad
[A,B]
\;\longleftrightarrow\;
i\hbar\,\{A,B\}_{\mathrm{PB}}.
\end{equation}
The second version is the one most often quoted.

The lecture ends with a quick check of energy conservation, but it is careful to state only what has actually been shown. In the time-translation-invariant case, where \(H\) has no explicit time dependence, we may set \(L=H\) in the expectation-value equation:
\begin{equation}
\frac{d}{dt}\langle H\rangle=\frac{i}{\hbar}\langle [H,H]\rangle.
\end{equation}
But any operator commutes with itself, so
\begin{equation}
[H,H]=0,
\end{equation}
and therefore
\begin{equation}
\frac{d}{dt}\langle H\rangle=0.
\end{equation}
This is the weak sense of energy conservation established in the lecture: the \emph{average} energy is constant in time. If \(H\) depends explicitly on time, an extra term appears and this simple conservation law fails, exactly as Susskind says it does in classical mechanics.

And here the lecture deliberately stops. The postulates used so far do not yet explain collapse. Between measurements the state evolves unitarily. A measurement is another matter, because the measurement process involves entangling the system with an apparatus. Until entanglement has been discussed, the notes can do no better than the lecture itself: they must postpone the derivation and temporarily adopt the language of collapse.

\section{Summary}

The lecture unfolds in a deliberate order. We first rebuild the measurement postulates: observables are represented by Hermitian operators, their eigenvalues are the possible outcomes of experiments, and complete physical distinguishability is encoded as orthogonality. Born's rule then turns overlaps into probabilities, with degeneracy handled by choosing an orthonormal basis in the degenerate subspace and adding the corresponding probabilities.

Only then do we ask the next question, the classical question: how do states change with time? The answer is a common linear time-development operator \(U(t)\), constrained by reversibility to preserve inner products and therefore to be unitary. Examining \(U(\epsilon)\) for very small \(\epsilon\) produces a Hermitian generator \(H\), the Hamiltonian, and from that the generalized time-dependent Schr\"odinger equation. Once the state-vector dynamics are in hand, the lecture turns to expectation values, derives the commutator equation for their evolution, compares it with the classical Poisson-bracket equation, and checks that the average energy is conserved when \(H\) has no explicit time dependence. The measurement problem is left explicitly unresolved. It will return only after entanglement is on the table.